\subsubsection*{Summing the local bounds}
\hypertarget{source-numerator-sums}{}

A bound for a congruence ideal must first be turned into a bound for the corresponding integral module. This conversion keeps the rank of the retained eigenspace in the sum over good characters.

\begin{established}[Hecke congruence ideal and normalized length]\label{hyp:eigencharacter-lattices}
First consider the passage from a Hecke congruence ideal
to the congruence length of its module. Work over one
coefficient DVR \(O\), normalize \(v(\ell)=1\), and write
\(Y\) for the theta eigensystems and \(Z\subseteq Y\) for
the good ones. Let \(T\) be the Hecke algebra inside the
algebra of functions on the generic eigensystems. Write \(I_Y=\{f\in T:f|_{Y^c}=0\}\). For each \(\pi\in Z\) its evaluation and normalized colength satisfy
\[
 0\ne I_Y(\pi)=(f_\pi(\pi))\subseteq O,\qquad
 v_{cc}(\pi)=v(f_\pi(\pi)),\qquad f_\pi\in I_Y .
\]
\end{established}
\begin{proof}
For each generic character \(y\ne\pi\), choose \(t_y\in T\) with \(t_y(\pi)\ne t_y(y)\). The product \(\prod_{y\ne\pi}(t_y-t_y(y))\) vanishes away from \(\pi\) and has nonzero evaluation there. The scalars \(t_y(y)\) belong to the splitting DVR \(O\). Hence the image of the supported ideal is a nonzero ideal of \(O\). Choose a generator and lift it to a supported element of \(T\).
\end{proof}

\begin{lemma}[From eigencharacter valuations to module colength]\label{lem:eigencharacter-colength}
With the ideal and length of Lemma~\ref{hyp:eigencharacter-lattices},
let \(M\) be the Hecke-stable lattice and \(V_Z\) its good generic
subspace. Put \(U=p_ZM\), \(B=M\cap V_Z\)
and let \(a\) be the integral type operator of
Proposition~\ref{prop:cut-selector}. Then the normalized
colength of \(B\) in \(U\) is at most
\[
 \sum_{\pi\in Z}\dim(V_\pi)\bigl(v_{cc}(\pi)+v_\ell(h)\bigr).
\]
\end{lemma}

\begin{proof}
The type operator \(a\) acts by \(h\) on the good
eigenspaces and by zero on the remaining eigenspaces
of \(Y\). The product \(af_\pi\) kills every eigenspace
outside \(Z\). It also preserves the integral module,
so it maps \(U=p_ZM\) into \(B=M\cap V_Z\), by applying
it to a lift in \(M\). The products for the different
\(\pi\) may all act on multiple good eigenspaces.
Order those eigenspaces and filter \(U\) and \(B\) by
intersection with their successive partial sums.

For
the step adding \(V_\pi\), the graded lattice of \(U\)
is the projection of the intersection at this step
onto \(V_\pi\), and likewise for \(B\). Multiplication
by \(h f_\pi(\pi)\) takes the first graded lattice into
the second, since \(af_\pi\) preserves the partial sum.
For the step adding \(V_\pi\), the DVR index formula therefore bounds the normalized graded length by
\[
 \dim V_\pi\bigl(v_{cc}(\pi)+v_\ell(h)\bigr).
\]
With \(U_i,B_i\) the intersections with the first \(i\) eigenspaces,
the successive index is read in the projection to the new eigenspace.
Equivalently
\[
 \operatorname{len}(U/B)=
 \sum_i\operatorname{len}\bigl((U_i/U_{i-1})/(B_i/B_{i-1})\bigr)
 \le\sum_{\pi}\dim V_\pi\bigl(v_{cc}(\pi)+v_\ell(h)\bigr).
\]
This proves the local comparison, with the Hecke
algebra and the theta set of the cut when it is
applied to the cut.
\end{proof}

\begin{proposition}[Weighted local return on the cut]\label{prop:weighted-local-return}
Fix a lift on the selected point family and use its good
multiplicities from Proposition~\ref{prop:cut-multiplicity}.
The numerator of Definition~\ref{def:weighted-congruence-numerator} satisfies \eqref{eq:local-division-star} of Section~\ref{subsec:division-length} on the cut, with the allowed additional term \((\log x)^2N_c\) and a threshold uniform over the selected points. In the notation of the local table its contribution at an ordinary localization \(i\), away from the two type primes, satisfies
\[
 \eta_i\le \sum_{\pi\in Z_i^+}m_{i,\pi}
               \bigl(v_{cc}(i,\pi)+v_\ell(h)\bigr).
\]
Here \(v_{cc}(i,\pi)\) is the local congruence bound, \(h\) the order of the finite row group and \(m_{i,\pi}\) the good multiplicity. The separate loss \(d_j\log h\) from returning to the original Hecke module is absorbed at this scale.
\end{proposition}

\begin{proof}
\proofstep{Weight each localization by its normalized local degree and good multiplicities.}
Apply Lemma~\ref{lem:eigencharacter-colength} on the cut at
a single selected point. Fix a lift \(j\) and write \(d_j\) for the rank after
the small-level rank identification. Let \(\mathcal T\)
be the finite set of rational primes below the twisted
coefficient places in its support. The type and auxiliary prime cases use
Propositions~\ref{prop:type-prime-row-length}
and~\ref{prop:auxiliary-prime-row-length}. The local length
comparison at a prime \(\ell\) other than
\(g_0,g_1\), is the bound
\[
 \eta_i\le \sum_{\pi\in Z_i^+}
      m_{i,\pi}\bigl(v_{cc}(i,\pi)+v_\ell(h)\bigr).
\]
Here \(\eta_i\) is the normalized congruence length of the
cut localization, \(m_{i,\pi}\) is its multiplicity, and
\(v_{cc}\) is the local congruence bound attached to the
Galois family at \(\pi\). The factor \(h\) belongs to this
comparison on the cut itself. Returning to the original
module incurs the separate loss proved above.

Lemma~\ref{lem:marked-row-column-lattices} gives the required integral presentation at each good character of this cut. Proposition~\ref{prop:marked-signature} supplies its complement estimate. The resulting local table has
\[
 v_{cc}(i,\pi)\le
 \nu_\ell+2\ell^{-2}+\mathcal L_\ell+V_\ell+v_\ell(h_E),
\]
where \(\nu_\ell\) and \(\mathcal L_\ell\) are independent
of \(i,\pi\) for fixed \(\ell\) at this point. The table
incorporates the uniform bound on \(s_\pi\) separately
from the split row estimate. At each twisted \(w\mid\ell\), \eqref{eq:retained-multiplicity}
gives
\[
 \eta_{j,w}=\sum_{i\text{ over }(j,w)}\eta_i
 \le d_j\bigl(\nu_\ell+2\ell^{-2}+\mathcal L_\ell+V_\ell
                 +v_\ell(h_E)+v_\ell(h)\bigr).
\]
All terms in this upper bound are nonnegative and independent of
\(w\). Multiply by \(e_w\log\mathrm Nw/[K:\mathbb Q]\) and sum
over the chosen places. The factor from their weights is at most
\(\log\ell\), by the residue-degree formula on all places above
\(\ell\). At the two type primes the
paired-character estimate bounds the weighted lengths
directly by
\(d_j\log\ell(\mathcal L_\ell+2\ell^{-2}+V_\ell+v_\ell(h_E))\).
Up to nonnegative terms the same summation applies there.

The class number term contributes at most \(\log h_E\),
since
\[
 \sum_{\ell\in\mathcal T}v_\ell(h_E)\log\ell=\log\!\left(\prod_{\ell\in\mathcal T}\ell^{v_\ell(h_E)}\right)\le\log h_E.
\]
The factor \(h\) costs at most \(\log h\)
by the same argument. Also
\(\sum_{\ell\in\mathcal T}\ell^{-2}\log\ell\) is absolutely
bounded by the series over the integers. The strict
cohomology terms contribute at most the global budget
\[
 V=\sum_\ell V_\ell\log\ell
\]
once, since \(\ell\) is
summed only after the multiplicity and degree estimates.
This bound is independent of the number of localizations
and the coefficient-field degree.

The scalar cyclotomic term has its own bound for the weighted sum over primes.
Put \(a=tr\), which is odd since \(t\) is odd, \(k,m\)
are even, and \(r=k+m-1\). The cyclotomic minimum \(\nu_\ell\) in \eqref{eq:signature-cyclotomic-minimum}
is taken over the absolute Galois group of this same field \(E\). If \(\ell\mid Q_S\), the image
contains \(1+\ell^{v_\ell(Q_S)}\) since \(E\) lies in
the cyclotomic field of conductor \(Q_Sqd\). The
binomial argument in Section~\ref{subsec:split-local-length}
gives
\[
 \nu_\ell\le v_\ell(Q_S)+v_\ell(a).
\]

The depth at \(2\) is at least two. If
\(\ell\nmid Q_Sqd\), the cyclotomic image is full by
the Chinese remainder theorem. Then \(\ell\) is odd
and \(-1\) gives \(\nu_\ell=0\).
At \(q,d\), the host Artin action of the factor at
that prime has order dividing two. Thus the cyclotomic
image modulo that prime contains the subgroup of
squares, of order \((\ell-1)/2\). This order exceeds
\(a\) by the auxiliary prime choice, so a generator raised
to \(a\) differs from \(1\). This again gives
\(\nu_\ell=0\). Therefore
\[
 \sum_{\ell\in\mathcal T}\nu_\ell\log\ell
 \le \log Q_S + \log(tr)=O(\log x).
\]
Here \(Q_S\le4 f_\psi n\le n^2\) eventually by
\eqref{eq:host-character-compositum} and \eqref{eq:dyadic-coordinate-orders}.
The target is fixed before the scale and the ranks grow polynomially.

The multiplicity estimate and the local table join at
one common threshold in \(x\). The arguments of
Section~\ref{subsec:special-prime-lengths}
give the type and ramified entries. Section~\ref{subsec:ordinary-split-length}
proves the ordinary split entry and the implication from a local
inert band to the indicator in the sum. The cyclic row estimates
in Section~\ref{subsec:inert-rows} supply the long-orbit and
the two remaining inert entries. The comparison above the
cutoff in Section~\ref{subsec:outside-windows} uses the
integral local classification on the same jump windows. The table
assigns \(\mathcal L_\ell\) in the following ranges.

\begin{itemize}
\item At each type prime \(g_i\) the table entry and sizes satisfy
\[
 \mathcal L_{g_i}\le Cx^C\eta_c,\qquad \log\eta_c\le x^{\epsilon/4},\qquad \log g_i=O(z).
\]

\item At each dyadic or odd conductor prime and at three,
Proposition~\ref{prop:small-split-loss} bounds the cost by \(C fmb\log x+Cx^Cz^2\). The weighted
sum is \(O(fmbz^2)\) by the size comparison of
Section~\ref{subsec:split-local-length}.
\item At \(q,d\), the cost is at most \(2C fmb\).
The weighted sum is \(O(fmbz)\) by the conductor of the auxiliary field
bound on the chosen prime pair.
\item The remaining table uses a positive cutoff
\(\ell_*\). The summation below gives \(O(fmbz^2)\)
for the boundary, split, inert and short-band costs
weighted by \(\log\ell\) at primes at most \(\ell_*\).
\item Above \(\ell_*\) the remainder is supported on
the exceptional set. With
\[
 U=x^{C+2}\exp(x^{\epsilon/4}),
\]
its cardinality is at most \(388U^8\). The weighted cost
per exceptional prime is at most \(3000(C+1)^2U^{10}\),
giving at most \(1164000(C+1)^2U^{18}\) in total.
\end{itemize}

\proofstep{Sum the same table before absorbing the exceptional-prime cost.}
The table is now weighted by the local degrees and actual multiplicities. Summing those entries first keeps the exceptional primes and the scalar cohomology on the same retained family.

For fixed \(C\) and \(\epsilon>0\),
\[
 \log(U^{18}/b)=18(C+2)\log x+18x^{\epsilon/4}-x^\epsilon
 \longrightarrow -\infty .
\]
In particular the exceptional sum and the two type-prime
costs are \(o(b)\) uniformly. The same comparison absorbs
any fixed polynomial in \(x\). The bound on \(\log h_E\)
can be obtained separately. All character conductors of the abelian
host divide \(N\), so the conductor--discriminant formula
gives \(\log|\Delta_E|\le [E:\mathbb Q]\log N\).
Minkowski's bound supplies an ideal in every class of
norm at most \(2^{[E:\mathbb Q]}|\Delta_E|^{1/2}\).
\cite[Theorems 3.27 and 4.3 respectively]{milne-cft,milne-ant}
With \(n_E=[E:\mathbb Q]\) and
\(D=2^{n_E}|\Delta_E|^{1/2}\), count these representatives by
\[
 h_E\le\sum_{\mathrm N\mathfrak a\le D}1
  \le D^2\sum_{\mathfrak a}(\mathrm N\mathfrak a)^{-2}
  =D^2\zeta_E(2)\le D^2\zeta(2)^{n_E}.
\]
The last inequality is Euler factor by Euler factor, since there
are at most \(n_E\) prime ideals above each rational prime and
their residue degrees are at least one.

Taking logarithms and
using \(\log|\Delta_E|\le n_E\log N\) gives
\(\log h_E=O(n_E(1+\log N))\), which is polynomial in \(x\)
under the selected size bounds. With the local table,
these estimates give the cut numerator bound (*) with the
allowed additional
nonnegative term \((\log x)^2 N_c\). The loss \(d_j\log h\)
from the cut to the original level can be absorbed
at this same scale.

The local inputs are assembled on the cut family
before applying the weighted sum. The estimate uses
the complete local table and the rank and period
identification at the chosen points. \label{cor:numerator-sums-1}
\end{proof}

\subsubsection*{Summation below the prime cutoff}
\label{subsec:prime-cost-sums}
\hypertarget{source-prime-cost-sums}{}

The remaining sum below the cutoff contains contributions from both split and inert places as well as boundary terms. Their individual bounds have to be summed with the prime weights before the total can be compared to the allowed numerator. Write \(z=x^\epsilon\), \(b=\exp z\), and
\(H=b+S_c+\eta_c+t+2\). Let \(k\) be the ambient rank parameter,
and put \(N(m)=k+m(t+1)\).
Write \(\delta\) for the density of the selected indices with the larger weights.
\begin{established}[Size inequalities for the local table]\label{hyp:small-prime-table-ranges}
The cutoff and the secondary threshold are
\[
 \ell_*=\frac{mb}{n t^2(N(m)+t)z^2+k n^2(S_c+t+z)},
 \qquad A=\frac{kz}{m\delta}.
\]
Write \(X'=\exp(x^{\epsilon/3})\), \(b'=xb/X'\), and
\(\omega=b/(nX')\). Let \(\eta_3=\eta_c+t+2\), let \(w_P\) be
the upper place-weight size plus \((t+1)/2\), and put
\(G_P=w_P+S_c+\eta_3\). For the selected ranks and weights, beyond a threshold depending on the fixed parameters,
\[
  \ell_*\le b/z,\quad z\le A,\quad r/A\le 2m/z,\quad(3+\log A)^2\le z,\quad H\le2b,\quad z\le b'.
\]
\[
 \begin{gathered}
 \eta_3\le4\eta_c,\qquad w_P+\eta_3\le9b',\qquad G_P\le10b',\\
 r b'A\log x\le mb,\qquad r\eta_3\le4m\omega,\qquad
 C(\eta_3+S_c+z)\le\omega/8 .
 \end{gathered}
\]
\end{established}
\begin{proof}
Lemma~\ref{lem:window-parameter-choice} fixes the ranks before the scale and the weights of the retained positions before choosing the place. In this
calculation \(x^{-\epsilon}\le\delta\le1\), \(r\le2k\),
\(R\le N(m)+t\), \(m\le k\), and the ranks are at most a fixed
power of \(x\). In particular \(\ell_*\le b/z\), and
\[
 z\le A,\qquad r/A\le 2m/z,\qquad
 (3+\log A)^2\le z
\]
beyond one threshold. The last inequality follows from
\(A\le k z^2\) and \(z/(\log x)^2\to\infty\).
The size inputs also give \(H\le2b\).

At the smaller places the two possible centers satisfy
\[
 \omega\le 5b'/4,\qquad b_v+\lfloor b'/x\rfloor\le5b'/4+b'/x .
\]
Here the bound for \(b_v\) is from Section~\ref{subsec:exceptional-range}. The widths \(S_c,\eta_c,t+2\) are \(o(b')\), and \(t+2\le3\eta_c\). Including these widths and the half shift gives \(w_P+\eta_3\le9b'\) and \(G_P\le10b'\) uniformly.

For the fourth inequality the selected reserve bounds
\(r(b'+S_c)A\log x\) by \(\delta mb\). The fifth follows from
the reserve \(N(m)\eta_c\le\delta m\omega\) and \(r\le N(m)\).
For the last inequality, the sizes of \(\eta_c,S_c,z\) are
\(O(\omega/\sqrt x)\), with constants fixed before taking \(x\)
large. For fixed \(\epsilon>0\) this follows also by comparing
their bounds, \(x^{O(1)}\exp(x^{\epsilon/4})\), with \(\omega\).
We have \(b'\ge z\) at large \(x\).
\end{proof}

\begin{proposition}[Summation up to the small-prime cutoff]\label{prop:small-prime-cost-total}
Under the cutoff and size inequalities of Lemma~\ref{hyp:small-prime-table-ranges},
let \(c_\ell\) be the sum of the local boundary, split and inert charges in the table at \(\ell\), before multiplication by \(\log\ell\). Then
\[
 \sum_{\ell\in P}(\log\ell)c_\ell\ll_C fmbz^2\qquad\bigl(P\subseteq\{\ell:5\le\ell\le\ell_*\}\text{ finite}\bigr).
\] The threshold
is uniform in the retained point and the subset of primes.
\end{proposition}

\begin{proof}
For a set of integers \(2\le a\le L\), summing over all such
integers gives
\[
 \sum \log a\le L\log L,\qquad
 \sum \frac{\log a}{a}\le (1+\log L)^2.
\]
The second bound uses the harmonic sum and \(\log a\le\log L\).
Replacing \(a\) by \(a-1\) in the denominator costs at most a
factor two. These bounds apply to any subset of the indicated primes.
The boundary cost \(Cfn t^2 R z^2\) consequently contributes at
most \(C fmbz\), by the first term of the cutoff denominator.
\proofstep{Split places.}
The split cost is \((2fm+f)e_P(\ell)\), where
\(e_P(\ell)=C(1+H/\ell+\log_\ell H)\). The other uses of \(e_P\)
have coefficient at most \(2Cfm\). After multiplication by
\(\log\ell\), the three terms in \(e_P\) sum to
\[
 O_C\!\left(\ell_*\log\ell_*+
             H(1+\log\ell_*)^2+\ell_*\log H\right)
       =O_C(bz^2).
\]

\proofstep{Inert places below the secondary threshold.}
For inert \(\ell<A\) of order exceeding \(4L_w\) in the
host's real cyclotomic quotient, the cost is \(4Cfr\) times
\[
 J_\ell=\frac{H}{(\ell-1)^{j_\ell+1}}
         +(G_P+1)(j_\ell+1)+z,\qquad
 j_\ell=\max(1,\lceil\log_{\ell-1}A\rceil).
\]
For \(\ell\ge5\), \((\ell-1)^{j_\ell}\ge A\) and
\(j_\ell+1\le3+\log A\). The logarithm-weighted sum of the
first term, including \(r\), is
\(O(rH(1+\log A)^2/A)=O(mb)\).
For the other terms it is
\(O(rb'A(3+\log A)^2)=O(mbz)\). This uses the reserve for
\(rb'A\) and \(G_P+1=O(b')\).

The other inert primes have a contribution \(4CfrH/\ell^2\).
At \(\ell\ge A\) use
\[
 \sum_{a\ge A}(\log a)/a^2=O((1+\log A)/A)
\]
by integral comparison, and \(r/A\le2m/z\). At \(\ell<A\)
the order is at most \(4L_w\); here the ordinary tail bound at
\(A\) gives no control. Put \(R_0=n/z\). The parameters of the auxiliary construction
give \(3\le R_0\), \(k/m\le R_0^2\), \(n\le o2^e\) and
\(2oL_w\le z\), where \(e\) is the host's dyadic exponent
and \(o\) its fixed cyclic order parameter. Thus
\[
 2^e/L_w\ge2R_0.
\]

The quotient
\((\mathbb Z/2^{e+2}\mathbb Z)^\times/\{\pm1\}\) has order \(2^e\),
so every order \(h\) on it is a power of two. By
Lemma~\ref{lem:short-real-dyadic-order} an integer of order \(h\) is
\(\pm1\pmod{2^{e+2-v_2(h)}}\).

Consequently \(h\le4L_w\) forces the original integer to
be \(\pm1\bmod M\), where
\[
 M=2^{e+2-\lfloor\log_2(4L_w)\rfloor}
       \ \ge\ 2^e/L_w\ \ge\ 2R_0 .
\]
The same congruence holds if \(a=1\). There is no term at
\(1\) in the sum of interest. All the other positive integers
in these two residue classes are \(jM\pm1\) with \(j\ge1\).
They are between \(jM/2\) and \(2jM\), so
\[
 \sum_{\substack{a\ge2\\a\equiv \pm1\bmod M}}
            \frac{\log a}{a^2}
   \le \frac8{M^2}\sum_{j\ge1}\frac{\log(2jM)}{j^2}
   =O\!\left(\frac{1+\log M}{M^2}\right).
\]
The function \((1+\log u)/u^2\) is decreasing for \(u\ge1\).
Multiplication by \(r\le2mR_0^2\) therefore gives
\(O(m(1+\log R_0))\). This bound accounts for the small
inert primes using only the dyadic quotient, with no estimate
for primes in arithmetic progressions.

The remaining inert terms cost \(4Cfr\) times
\[
  \frac{w_P+\eta_3}{\ell}+S_c+z+\eta_3 I_\ell .
\]
Here \(I_\ell\) is the indicator that some \(s\) between
\(\omega\) and the upper place-weight size before the
\((t+1)/2\) shift satisfies
\[
 |\ell s-b|\le
   C\bigl(\ell(\eta_3+S_c+z)+b/\ell+s\bigr).
\]
For \(\ell\ge2C\) the right side is at most
\(\ell s/8+b/2+\ell s/2\), by the selected reserves.
It follows that \(\omega\le s\le4b/\ell\) in this band.
Hence \(r\eta_3 I_\ell\le16mb/\ell\).

At the finitely
many \(\ell<2C\) we use \(r\eta_3\le4m\omega\le4mb\).
Also \(r(w_P+\eta_3)\le9rb'\le9mb\), and the second
term of the cutoff denominator gives
\[
 \ell_*r(S_c+z)\le2mb.
\]
The first and band terms
are summed using \(\sum (\log\ell)/\ell\), and the
constant term using \(\sum\log\ell\). Their total
contribution is \(O_C(fmbz^2)\).

The estimates above apply together to any finite subset of
the primes \(5\le\ell\le\ell_*\). They give the desired
\(O_C(fmbz^2)\) bound for the numerical table used in the
weighted export. The threshold is common to all the retained points with the specified size inputs.
The sum is therefore available for the actual row once each localization satisfies the corresponding local Galois estimate. \label{cor:prime-cost-sums-1}
\end{proof}

\subsubsection*{The exceptional range}
\label{subsec:exceptional-range}

Above \(\ell_*\), only the exceptional jump windows require this additional sum. We count the primes they can contain and bound the total weighted cost uniformly in the window data. Write
\[
 \mathcal D=mb/\ell_*,\qquad
 U=x^{C+2}\exp(x^{\epsilon/4}),\qquad
 s_+=\lfloor2b/\ell_*\rfloor .
\]
The selected parameters give \(U\ge1\),
\(n,k,t,m,z,\eta_c\le U\), \(f,r\le2U\),
\(R,N(m)+t\le4U^2\) and \(S_c\le4U\).
Fixed polynomial rank bounds suffice after taking
\(C\) large enough. The two summands of \(\mathcal D\)
are at most \(4U^7\) and \(6U^4\) respectively, hence
\[
 \mathcal D\le10U^7.
\]
Since \(U^{14}=o(b)\), this
gives \(\ell_*^2>8b\) beyond a common threshold.
Also \(\ell_*>A\), since \(A\le k z^2\) grows
polynomially in \(x\).

The centers \(W\) of the exceptional windows are
\(b,\omega\), and \(b_v+\lfloor b'/x\rfloor\) for the
top places \(v\). Their number is at most \(2n+2\),
using the bound \(2n\) for the number of top places.
The weight construction gives
\(\omega\le b_v\le 5b'/4\), so each center lies in
\([\omega,2b]\) for large \(x\). The window about \(W\)
contains the integers in
\[
[W-\eta_c-2s_+-3,\,
  W+S_c+2\eta_c+2s_++4].
\]
All these integers are positive and less than \(3b\).
Indeed \(s_+\le2\mathcal D/m\) and
\(U^7=o(\omega)\) since \(n\) is polynomial,
\(\omega=b/(nX')\) and \(\log X'=x^{\epsilon/3}\).
This controls both endpoint displacements.

The exceptional set consists of primes \(\ell>\ell_*\)
dividing one of these integers. Two distinct such
divisors would have product greater than
\(\ell_*^2>8b\), so each integer contributes at most
one prime. An interval of length \(D\) contains at
most \(D+1\) integers. Thus the set has size at most
\[
 (2n+2)(S_c+3\eta_c+4s_++8)
 \le(4U)(97U^7)=388U^8.
\]
Each prime in the set is less than \(3b\), so
\(\log\ell\le3U\).
\label{cor:prime-cost-sums-2}

For these primes the boundary cost is
\(Cfn t^2 R z^2\le8C U^8\). The envelopes for the
pointwise row costs are the split cost and the second
inert cost, since \(\ell>\ell_*>A\). For \(\ell>7\),
\[
\begin{split}
 H/\ell^2 &\le H/\ell \le 2\mathcal D\le20U^7,\\
 e_P(\ell)&=C(1+H/\ell+\log_\ell H)\le23C U^7.
\end{split}
\]
For the flat part of the inert cost, the numerical
expression is
\((w_P+\eta_3)/\ell+S_c+z+\eta_3 I_\ell\),
where \(I_\ell\in\{0,1\}\) is the band indicator.
Here \(b'<\ell_*\) at the common threshold, by
\(\mathcal D x/X'\to0\), and \(w_P+\eta_3\le9b'\).
Thus the flat part is at most
\(9+4U+U+4U\le18U^7\).

The split envelope is at most \(138C U^9\).
The inert envelope is at most
\(608C U^9+92C^2U^9\), using
\[
 4Cfr(20+18)U^7+2Cfm(23C U^7)
       \le (608C+92C^2)U^9.
\]
Adding the boundary cost and multiplying by
\(\log\ell\) gives the bound
\(3000(C+1)^2 U^{10}\) per prime.
This proves the exceptional contribution to the
weighted sum from the pointwise row comparisons.
The complementary large-prime range uses the
congruence estimate with zero additional length
outside these windows.
\label{cor:prime-cost-sums-3}

